\section{Discussion and Limitations}
\label{sec:discussion}

%% Summary of sections 3-5; no new measurements.
\begin{table}[t]
\centering
\small
\setlength{\tabcolsep}{3pt}
\begin{tabular}{@{}p{0.15\columnwidth}p{0.42\columnwidth}p{0.37\columnwidth}@{}}
\toprule
Layer & Provides & Does not provide \\
\midrule
Skills & measured lessons that persist across sessions; control of what each sub-agent sees
       & correctness: compliance is not proof \\
PBT    & calibrated, replayable verdicts; detection equal to a reference
       & anything beyond tested inputs, or a complete bundle \\
Formal & a model of the kernel proven to meet a spec that is pinned between the reference and known wrong kernels
       & device IEEE behaviour, compiler; completeness beyond the wrong kernels held \\
\bottomrule
\end{tabular}
\caption{What each layer contributes to a kernel the loop keeps. Of the formal row,
  only the checking core is built.}
\label{tab:layers}
\end{table}

Table~\ref{tab:layers} summarizes the division of labour. The layers fail independently,
and each covers a failure the others cannot see. Skills cannot make a gate complete. A
calibrated gate can still be incomplete (\S\ref{sec:incomplete}). A proof says nothing
about the code that actually runs unless it is checked against that code's executions,
which is why C2--C6 run on PBT's cases. The loop improves only when all three hold.

The results also sharpen what a declarative oracle is for. Where a cheap reference
exists, laws tie it on detection, add no calibration advantage over a well-normalized
ratio, and, as drafted, admit wrong kernels that the reference rejects. Their value lies
where no reference exists, as in translation to a new backend, and only once their
completeness has been checked. The same holds for an agent-written formal spec, which is
why the certificate design measures spec completeness (C3) instead of trusting it.

\paragraph{Threats to validity.}
Device results come from one GPU model, float32 only and seed 42. The mutation corpora
are small, at most 15 mutants each and all on softmax, and the fault class of a mutant is
the one its prompt asked for, not one that was verified. Absolute detection rates are
deflated by the ladder's single-column rungs (\ladderCeiling{} ceiling). Arm-versus-arm
comparisons are unaffected. The saboteurs and the families of Table~\ref{tab:gap} were
chosen by authors who knew the bundles. They show that the gate \emph{can} fail and are
not detection rates. The six rolled, reversed, sorted and negated rows come from a
re-run script not yet recorded with the project's measurements. The training results in
\S\ref{sec:incomplete} use one seed each. The gate counts abstention as a pass, which our
matrix guards only for its own rows. Of the formal layer, only the checking core exists.
It is tested against a fake adapter, and the guarantee that an adapter's executable
readings come from the proved definitions is still that adapter's obligation.
Proposition~\ref{prop:softmax}(a) and (b) are mechanized over an abstract ordered field and
(c) is proven on paper; all hold over $\R$ and leave the floating-point forms open. C3 is a lower bound, since a spec that rejects every wrong
kernel we hold may admit one nobody has written yet. Finally,
this area produces a closely competing preprint roughly every month, so the related work
must be re-checked before any submission.

\section{Conclusion}

For an agent that optimizes kernels, the correctness gate is half the objective. We
showed that its calibration decides whether the loop can start, that its completeness
decides what the loop converges to, and that a gate built from plausible laws can admit a
wrong kernel \fillSpeedup$\times$ faster than the real one. Skills carry those lessons
forward and keep each experiment's agents honest. PBT makes the gate cheap and fair to
compare. Formal proof, written by the agent but checked against PBT's own executions
and issued only after the gate accepts a kernel, is what can turn
``passes the gate'' into a guarantee.
